\documentclass[11pt,a4paper]{article}
% ============================================================
%  פרמבולה משותפת לחומרי הבחינה — מודולים 16–18
%  מהדרים עם:  lualatex <file>.tex
% ============================================================
\usepackage{fontspec}
\directlua{
  luaotfload.add_fallback("hbfb", {"DejaVu Serif:mode=node;", "DejaVu Sans:mode=node;"})
  luaotfload.add_fallback("ltfb", {"DejaVu Sans:mode=node;"})
}
\usepackage[bidi=basic]{babel}
\babelprovide[import, main]{hebrew}
\babelprovide[import]{english}
\babelfont{rm}[Ligatures=TeX, RawFeature={fallback=ltfb}]{DejaVu Serif}
\babelfont[hebrew]{rm}[Script=Hebrew, Ligatures=TeX, RawFeature={fallback=hbfb}]{Noto Serif Hebrew}
\babelfont{tt}[Scale=0.84, RawFeature={fallback=ltfb}]{DejaVu Sans Mono}

\usepackage[a4paper,top=2.1cm,bottom=2.0cm,left=1.9cm,right=1.9cm,headsep=0.5cm,footskip=1cm]{geometry}
\usepackage{booktabs, longtable, array, adjustbox, xcolor, enumitem, fancyhdr, titlesec,
            fancyvrb, needspace, multicol, lastpage, tikz, colortbl}
\usetikzlibrary{positioning, fit, backgrounds, matrix}

\definecolor{inkblue}{HTML}{14365A}
\definecolor{rulegray}{HTML}{8A97A3}
\definecolor{boxgray}{HTML}{F4F6F8}
\definecolor{codebg}{HTML}{F7F8FA}
\definecolor{warnred}{HTML}{9B2226}
\definecolor{okgreen}{HTML}{1B5E3A}
\definecolor{tipgold}{HTML}{8A6D1F}

\setlength{\parindent}{0pt}
\setlength{\parskip}{3pt plus 1pt}
\renewcommand{\arraystretch}{1.15}
\setlength{\columnsep}{18pt}
\setlength{\headheight}{13pt}

\pagestyle{fancy}\fancyhf{}
\renewcommand{\headrulewidth}{0.4pt}
\renewcommand{\footrulewidth}{0pt}
\fancyfoot[C]{\small עמוד \thepage{} מתוך \pageref{LastPage}}

% ---- כותרות ----
\titleformat{\section}{\needspace{3\baselineskip}\large\bfseries\color{inkblue}}{\thesection.\ }{0pt}{}
\titlespacing*{\section}{0pt}{15pt plus 3pt}{6pt}
\titleformat{\subsection}{\needspace{2.5\baselineskip}\normalsize\bfseries\color{inkblue}}{\thesubsection\ }{0pt}{}
\titlespacing*{\subsection}{0pt}{10pt plus 2pt}{4pt}

% ---- קו מפריד מתחת לכותרת חלק ----
\newcommand{\partheader}[3]{%
  \section*{#1}%
  \vspace{-4pt}{\color{rulegray}\hrule height 0.8pt}\vspace{4pt}%
  {\small\textbf{בסיס הנתונים:} #2 \hfill \textbf{ניקוד:} #3}\par\vspace{2pt}%
  {\color{rulegray}\hrule height 0.4pt}\vspace{6pt}}

\newcommand{\rulebar}{\vspace{-4pt}{\color{rulegray}\hrule height 0.8pt}\vspace{5pt}}

% ---- שאלה וסעיף ----
\newcounter{qnum}
\newcommand{\question}[2]{%
  \needspace{5\baselineskip}\refstepcounter{qnum}\par\vspace{8pt}%
  \noindent\textbf{\color{inkblue}שאלה \theqnum.}\ \textbf{#2}\hfill{\small(#1 נק')}\par\vspace{2pt}}
\newcommand{\subq}[2]{\needspace{3\baselineskip}\par\vspace{4pt}\noindent\textbf{#1.}~#2\par}
\newcommand{\subqp}[3]{\needspace{8\baselineskip}\par\vspace{5pt}\noindent%
  \textbf{#1.}~{\small\textbf{\color{inkblue}(#3 נק')}}~#2\par}
\newcommand{\pairitem}[1]{\needspace{2\baselineskip}\par\vspace{4pt}\noindent\hspace{1.4em}\textbf{#1}\par}

% ---- שורות תשובה מנוקדות ----
\newcommand{\answerlines}[1]{%
  \par\vspace{1pt}%
  \needspace{\dimexpr 6mm*#1 + 2mm\relax}%
  \begingroup\color{rulegray}%
  \count255=#1\relax
  \loop\ifnum\count255>0
    \noindent\babelsublr{\rule[-2pt]{0.6pt}{9pt}\rule[-2pt]{\dimexpr\linewidth-0.6pt}{0.3pt}}\par\vspace{5.4mm}%
    \advance\count255 by -1
  \repeat
  \endgroup\vspace{-2mm}}

% ---- שורת תשובה קצרה אחת ----
\newcommand{\blankline}[1]{\babelsublr{\rule[-2pt]{#1}{0.4pt}}}

% ---- קוד SQL, תמיד משמאל לימין ----
\DefineVerbatimEnvironment{sqlverb}{Verbatim}{fontsize=\small, frame=leftline, framerule=1.2pt,
  rulecolor=\color{rulegray}, framesep=7pt, xleftmargin=12pt, samepage=true}
\DefineVerbatimEnvironment{sqlbad}{Verbatim}{fontsize=\small, frame=leftline, framerule=1.2pt,
  rulecolor=\color{warnred}, framesep=7pt, xleftmargin=12pt, samepage=true}
\DefineVerbatimEnvironment{sqlgood}{Verbatim}{fontsize=\small, frame=leftline, framerule=1.2pt,
  rulecolor=\color{okgreen}, framesep=7pt, xleftmargin=12pt, samepage=true}

\newcommand{\lr}[1]{\babelsublr{\mbox{#1}}}
\newcommand{\sqlin}[1]{\babelsublr{\texttt{#1}}}

% ---- תיבה אפורה ----
\newsavebox{\panelbox}
\newenvironment{panel}{\par\addvspace{7pt}\begin{lrbox}{\panelbox}%
  \begin{minipage}{\dimexpr\linewidth-2\fboxsep-2\fboxrule}}%
  {\end{minipage}\end{lrbox}\noindent\fcolorbox{rulegray}{boxgray}{\usebox{\panelbox}}\par\addvspace{7pt}}

% ---- תיבה צבעונית עם כותרת ----
\newsavebox{\notebox}
\newenvironment{notepanel}[2]{\def\npcolor{#1}\par\addvspace{7pt}\begin{lrbox}{\notebox}%
  \begin{minipage}{\dimexpr\linewidth-2\fboxsep-2\fboxrule}%
  \textbf{\color{#1}#2}\par\vspace{2pt}}%
  {\end{minipage}\end{lrbox}\noindent\fcolorbox{\npcolor}{white}{\usebox{\notebox}}\par\addvspace{7pt}}

% ---- טבלת נתונים מלאה ----
\newcommand{\thd}[1]{\textbf{#1}}
\newenvironment{fulltable}[3]{%
  \needspace{5\baselineskip}\par\addvspace{7pt}%
  \noindent{\small הטבלה \babelsublr{\texttt{\textbf{#1}}} \textendash\ #2#3}\par\vspace{2pt}%
  \begin{otherlanguage}{english}\begin{adjustbox}{max width=\linewidth}}
  {\end{adjustbox}\end{otherlanguage}\par\addvspace{4pt}}

% ---- פלט של שאילתה ----
\newenvironment{outtable}[1]{%
  \needspace{4\baselineskip}\par\vspace{2pt}%
  \noindent{\small\textbf{הפלט}#1}\par\vspace{1pt}%
  \begin{otherlanguage}{english}\begin{adjustbox}{max width=\linewidth}}
  {\end{adjustbox}\end{otherlanguage}\par\addvspace{5pt}}

% ---- פלט ללא כותרת ----
\newenvironment{outplain}{%
  \par\vspace{2pt}\begin{otherlanguage}{english}\begin{adjustbox}{max width=\linewidth}}
  {\end{adjustbox}\end{otherlanguage}\par\addvspace{4pt}}

\newcommand{\nullv}{\textit{\small NULL}}
\providecommand{\tightlist}{\setlength{\itemsep}{0pt}\setlength{\parskip}{0pt}}


\fancyhead[R]{\small בסיסי נתונים ו\lr{-}\lr{SQL} \textbar{} בחינת מתכונת}
\fancyhead[L]{\small תרגול לקראת בחינת האמצע}

\begin{document}

\begin{center}
{\LARGE\bfseries\color{inkblue} בחינת מתכונת \textendash{} תרגול לפני הבחינה}\\[3pt]
{\large בית הקפה "אל-בלד" \ \textbullet\ \ \texttt{SELECT} \textbullet{} \texttt{WHERE} \textbullet{} \texttt{ORDER BY}}\\[6pt]
{\large\bfseries שעה אחת \ \textbar\ \ \lr{60} נקודות \ \textbar\ \ הפתרונות המלאים בסוף הדף}
\end{center}
\vspace{2pt}
{\color{inkblue}\hrule height 1.2pt}
\vspace{8pt}

\begin{panel}
\textbf{איך מתרגלים עם הדף הזה}
\begin{itemize}[leftmargin=1.3em, itemsep=2pt, topsep=3pt]
  \item הבחינה הזאת \textbf{קצרה וקלה יותר} מבחינת האמצע, ובנויה \textbf{באותו מבנה} בדיוק: ארבע שאלות, באותו סדר נושאים.
  \item \textbf{פתרו הכול לבד, על נייר, בשעה אחת} \textendash{} בלי מחשב ובלי להציץ. רק אחר כך פתחו את הפתרונות שבעמוד האחרון.
  \item מי שמסיים לפני הזמן: טענו את \sqlin{cafe.sql} ל-\lr{\texttt{onecompiler.com/sqlite}} והריצו את התשובות שלכם.
  \item \textbf{כל השאילתות כאן הן על טבלה אחת.} אין צורך לצרף טבלאות.
  \item ערך טקסט בגרש \textbf{בודד}, כינוי בגרש \textbf{כפול}, ונקודה-פסיק בסוף כל שאילתה.
\end{itemize}
\end{panel}

\vspace{6pt}
\section*{בסיס הנתונים: בית הקפה "אל-בלד"}
\vspace{-5pt}{\color{rulegray}\hrule height 0.8pt}\vspace{6pt}

בבית הקפה מועסקים \textbf{עובדים}, בתפריט יש \textbf{מנות}, וכל \textbf{הזמנה} היא מנה אחת שלקוח הזמין אצל עובד מסוים, בתאריך מסוים.

\vspace{2pt}
{\small\textbf{הערות לנתונים:} \lr{\texttt{HourPay}} הוא שכר העובד לשעה בש"ח, \lr{\texttt{DishPrice}} הוא מחיר המנה בש"ח, ו-\lr{\texttt{Guests}} הוא מספר הסועדים שישבו בשולחן. \textbf{לשני עובדים לא רשום טלפון.}}

\vspace{4pt}
\noindent\begin{minipage}[t]{0.37\linewidth}
\begin{fulltable}{Cities}{5 שורות}{}
\begin{tabular}{ll}
\toprule
\thd{CityCode} & \thd{CityName} \\
\midrule
1 & ירכא \\
2 & עכו \\
3 & כרמיאל \\
4 & נהריה \\
5 & שפרעם \\
\bottomrule
\end{tabular}
\end{fulltable}
\end{minipage}\hfill
\begin{minipage}[t]{0.59\linewidth}
\begin{fulltable}{Dishes}{8 שורות}{}
\begin{tabular}{llll}
\toprule
\thd{DishCode} & \thd{DishName} & \thd{DishKind} & \thd{DishPrice} \\
\midrule
10 & חומוס & מנה ראשונה & 28 \\
20 & סלט ירקות & מנה ראשונה & 32 \\
30 & שניצל & מנה עיקרית & 56 \\
40 & פסטה & מנה עיקרית & 48 \\
50 & קנאפה & קינוח & 30 \\
60 & עוגת גבינה & קינוח & 26 \\
70 & לימונדה & שתייה & 14 \\
80 & קפה & שתייה & 12 \\
\bottomrule
\end{tabular}
\end{fulltable}
\end{minipage}

\vspace{2pt}
\begin{fulltable}{Workers}{6 שורות}{}
\begin{tabular}{llllll}
\toprule
\thd{WorkerNum} & \thd{WorkerName} & \thd{WorkerCity} & \thd{StartDate} & \thd{HourPay} & \thd{WorkerPhone} \\
\midrule
1 & סאלי חדאד & 2 & 2022-03-01 & 42 & 050-1111111 \\
2 & יובל נחום & 1 & 2023-07-15 & 38 & \nullv \\
3 & אמיר טאהא & 2 & 2021-11-10 & 50 & 054-3333333 \\
4 & מיכל שרון & 4 & 2024-01-20 & 36 & 052-4444444 \\
5 & רים עבאס & 1 & 2020-05-05 & 55 & \nullv \\
6 & תום בראל & 5 & 2023-09-01 & 38 & 053-6666666 \\
\bottomrule
\end{tabular}
\end{fulltable}

\begin{fulltable}{Orders}{12 שורות}{}
\begin{tabular}{llllll}
\toprule
\thd{OrderNum} & \thd{CustName} & \thd{OrderDish} & \thd{OrderWorker} & \thd{OrderDate} & \thd{Guests} \\
\midrule
101 & דנה לוי & 30 & 1 & 2026-02-04 & 2 \\
102 & עומר חורי & 10 & 3 & 2026-02-18 & 4 \\
103 & נועה כץ & 40 & 2 & 2026-03-05 & 3 \\
104 & סמיר דיב & 30 & 1 & 2026-03-22 & 6 \\
105 & ליאת בר & 50 & 4 & 2026-04-09 & 2 \\
106 & ראמי סעיד & 20 & 3 & 2026-04-27 & 5 \\
107 & יעל שגיא & 30 & 5 & 2026-05-11 & 2 \\
108 & חנאן זועבי & 70 & 6 & 2026-05-30 & 8 \\
109 & איתי פלד & 40 & 2 & 2026-06-14 & 3 \\
110 & מאיה לוי & 60 & 4 & 2026-06-28 & 4 \\
111 & ואיל נסאר & 30 & 5 & 2026-07-07 & 2 \\
112 & שני אדרי & 80 & 6 & 2026-07-21 & 10 \\
\bottomrule
\end{tabular}
\end{fulltable}

\newpage
\question{10}{מבנה בסיס הנתונים}

\subqp{א}{כתבו את \textbf{המפתח הראשי} של כל אחת מהטבלאות.}{3}
\vspace{2pt}
\noindent\begin{tabular}{@{}p{0.20\linewidth}p{0.72\linewidth}@{}}
\lr{\texttt{Workers}} & \blankline{7cm} \\[7pt]
\lr{\texttt{Dishes}} & \blankline{7cm} \\[7pt]
\lr{\texttt{Orders}} & \blankline{7cm} \\[5pt]
\end{tabular}

\subqp{ב}{בטבלת \lr{\texttt{Orders}} יש \textbf{שני} מפתחות זרים. כתבו כל אחד, ואת הטבלה והשדה שאליהם הוא מצביע.}{4}
\vspace{2pt}
\noindent\begin{tabular}{@{}p{0.30\linewidth}p{0.30\linewidth}p{0.32\linewidth}@{}}
\toprule
\thd{המפתח הזר} & \thd{מצביע לטבלה} & \thd{לשדה} \\
\midrule
\blankline{3.4cm} & \blankline{3.4cm} & \blankline{3.4cm} \\[8pt]
\blankline{3.4cm} & \blankline{3.4cm} & \blankline{3.4cm} \\[4pt]
\bottomrule
\end{tabular}

\subqp{ג}{מהו \textbf{סוג הקשר} בין \lr{\texttt{Cities}} ל-\lr{\texttt{Workers}}? נמקו במשפט אחד.}{3}
\answerlines{3}

\vspace{6pt}
\question{15}{\lr{\texttt{SELECT}} \textendash{} מה מציגים}

\subqp{א}{הציגו את שם המנה ואת מחירה. תנו לעמודות כותרות בעברית: \textbf{שם המנה} ו-\textbf{מחיר}.}{5}
\answerlines{3}

\subqp{ב}{אילו \textbf{סוגי מנות} קיימים בתפריט? הציגו כל סוג \textbf{פעם אחת בלבד}.}{5}
\answerlines{2}

\subqp{ג}{בית הקפה נותן \textbf{הנחה של \lr{10\%}} על כל המנות. הציגו את שם המנה, את המחיר, ואת המחיר אחרי ההנחה תחת הכותרת \lr{\texttt{NewPrice}}.}{5}
\answerlines{3}

\newpage
\question{20}{\lr{\texttt{WHERE}} \textendash{} אילו שורות}

\subqp{א}{הציגו את שם המנה ואת מחירה, של המנות שמחירן \textbf{גבוה מ-\lr{30}} ש"ח.}{5}
\answerlines{3}

\subqp{ב}{הציגו את מספר ההזמנה, שם הלקוח והתאריך, של ההזמנות שהתקבלו \textbf{בין \lr{1.4.2026} ל-\lr{30.6.2026}} (כולל). השתמשו ב-\lr{\texttt{BETWEEN}}.}{5}
\answerlines{3}

\subqp{ג}{הציגו את מספר העובד ואת שמו, של העובדים ש\textbf{לא רשום להם טלפון}.}{5}
\answerlines{3}

\subqp{ד}{הציגו את מספר ההזמנה ואת שם הלקוח, של כל הלקוחות ש\textbf{שם המשפחה שלהם "לוי"}. שימו לב: בעמודה \lr{\texttt{CustName}} שמור השם המלא.}{5}
\answerlines{3}

\vspace{6pt}
\question{15}{\lr{\texttt{ORDER BY}} \textendash{} באיזה סדר}

\subqp{א}{הציגו את שם המנה, הסוג והמחיר של \textbf{כל} המנות, ממוינות לפי \textbf{סוג המנה} בסדר עולה, ובתוך כל סוג \textendash{} לפי \textbf{המחיר מהגבוה לנמוך}.}{5}
\answerlines{3}

\subqp{ב}{הציגו את \textbf{שלוש ההזמנות עם הכי הרבה סועדים}: מספר הזמנה, שם לקוח ומספר הסועדים.}{5}
\answerlines{3}

\subqp{ג}{הציגו את שם העובד ואת שכרו לשעה, ממוינים לפי \textbf{השכר מהגבוה לנמוך}. לשני עובדים אותו שכר \textendash{} ביניהם מיינו לפי \textbf{השם} בסדר האלף-בית.}{5}
\answerlines{3}

\vspace{10pt}
{\color{inkblue}\hrule height 1pt}
\vspace{4pt}
\begin{center}
{\small\textbf{סיימתם?} בדקו שכל שאילתה מסתיימת ב\textbf{נקודה-פסיק}, ושכל ערך טקסט כתוב ב\textbf{גרש בודד} \textendash{} ורק אז הפכו את הדף.}
\end{center}

\newpage
\begin{center}
\vspace*{2.5cm}
{\color{warnred}\hrule height 1.5pt}
\vspace{10pt}
{\Huge\bfseries\color{warnred} עצרו}\\[10pt]
{\LARGE\bfseries אל תמשיכו לפני שסיימתם}\\[12pt]
{\large בעמודים הבאים נמצאים \textbf{הפתרונות המלאים}.}\\[4pt]
{\large תרגול שבו מציצים בתשובה לפני שמנסים \textendash{} אינו תרגול.}
\vspace{10pt}
{\color{warnred}\hrule height 1.5pt}
\end{center}

\newpage
\partheader{פתרונות \textendash{} שאלה 1}{\texttt{cafe.sql}}{10 נקודות}

\subqp{א}{מפתחות ראשיים}{3}
\vspace{2pt}
\noindent\begin{tabular}{@{}p{0.22\linewidth}p{0.30\linewidth}p{0.42\linewidth}@{}}
\toprule
\thd{הטבלה} & \thd{המפתח הראשי} & \thd{למה} \\
\midrule
\lr{\texttt{Workers}} & \lr{\texttt{WorkerNum}} & המספר מזהה עובד אחד, ואינו חוזר \\
\lr{\texttt{Dishes}} & \lr{\texttt{DishCode}} & הקוד מזהה מנה אחת \\
\lr{\texttt{Orders}} & \lr{\texttt{OrderNum}} & מספר ההזמנה מזהה הזמנה אחת \\
\bottomrule
\end{tabular}

\subqp{ב}{מפתחות זרים בטבלת \lr{\texttt{Orders}}}{4}
\vspace{2pt}
\noindent\begin{tabular}{@{}p{0.26\linewidth}p{0.26\linewidth}p{0.24\linewidth}p{0.18\linewidth}@{}}
\toprule
\thd{המפתח הזר} & \thd{מצביע לטבלה} & \thd{לשדה} & \thd{הקשר} \\
\midrule
\lr{\texttt{OrderDish}} & \lr{\texttt{Dishes}} & \lr{\texttt{DishCode}} & \lr{1:N} \\
\lr{\texttt{OrderWorker}} & \lr{\texttt{Workers}} & \lr{\texttt{WorkerNum}} & \lr{1:N} \\
\bottomrule
\end{tabular}

\vspace{3pt}
{\small\lr{\texttt{CustName}} \textbf{אינו} מפתח זר \textendash{} הוא טקסט חופשי, ואין טבלה של לקוחות.}

\subqp{ג}{הקשר בין \lr{\texttt{Cities}} ל-\lr{\texttt{Workers}}}{3}

\vspace{2pt}
התשובה: \textbf{\lr{1:N}} (אחד לרבים).

\vspace{2pt}
\textbf{הנימוק:} בעיר אחת יכולים לעבוד כמה עובדים, אבל לכל עובד רשומה \textbf{עיר אחת} בלבד. המפתח הזר \lr{\texttt{WorkerCity}} יושב בטבלת \lr{\texttt{Workers}} \textendash{} כלומר בצד ה"רבים" \textendash{} והערכים בו חוזרים (ירכא ועכו מופיעות פעמיים כל אחת).

\begin{panel}
\textbf{בונוס למי שבדק:} לעיר \textbf{כרמיאל} (קוד \lr{3}) אין אף עובד. זה \textbf{תקין לגמרי} \textendash{} בקשר \lr{1:N}, לשורה בצד ה"אחד" מותר שלא יהיו לה שורות מתאימות בכלל.
\end{panel}

\vspace{4pt}
\partheader{פתרונות \textendash{} שאלה 2}{\texttt{cafe.sql}}{15 נקודות}

\subqp{א}{שם המנה ומחירה, עם כותרות בעברית}{5}
\begin{otherlanguage}{english}
\begin{sqlverb}
SELECT DishName AS "שם המנה", DishPrice AS "מחיר"
FROM   Dishes;
\end{sqlverb}
\end{otherlanguage}

\begin{outtable}{ \textendash{} \lr{8} שורות}
\begin{tabular}{ll}
\toprule
\thd{שם המנה} & \thd{מחיר} \\
\midrule
חומוס & 28 \\
סלט ירקות & 32 \\
שניצל & 56 \\
פסטה & 48 \\
קנאפה & 30 \\
עוגת גבינה & 26 \\
לימונדה & 14 \\
קפה & 12 \\
\bottomrule
\end{tabular}
\end{outtable}

{\small\textbf{שימו לב:} כותרת בעברית \textendash{} או כל כותרת עם רווח \textendash{} חייבת גרשיים \textbf{כפולים}.}

\subqp{ב}{סוגי המנות, כל סוג פעם אחת}{5}
\begin{otherlanguage}{english}
\begin{sqlverb}
SELECT DISTINCT DishKind FROM Dishes;
\end{sqlverb}
\end{otherlanguage}

\begin{outplain}
\begin{tabular}{l}
\toprule
\thd{DishKind} \\
\midrule
מנה ראשונה \\
מנה עיקרית \\
קינוח \\
שתייה \\
\bottomrule
\end{tabular}
\end{outplain}

{\small מ-\lr{8} שורות נשארו \lr{4}. בלי \lr{\texttt{DISTINCT}} היו חוזרות כל שמונה השורות, וכל סוג היה מופיע פעמיים.}

\subqp{ג}{מחיר אחרי הנחה של \lr{10\%}}{5}
\begin{otherlanguage}{english}
\begin{sqlverb}
SELECT DishName, DishPrice, DishPrice * 0.9 AS NewPrice
FROM   Dishes;
\end{sqlverb}
\end{otherlanguage}

\begin{outtable}{ \textendash{} \lr{8} שורות}
\begin{tabular}{lll}
\toprule
\thd{DishName} & \thd{DishPrice} & \thd{NewPrice} \\
\midrule
חומוס & 28 & 25.2 \\
סלט ירקות & 32 & 28.8 \\
שניצל & 56 & 50.4 \\
פסטה & 48 & 43.2 \\
קנאפה & 30 & 27.0 \\
עוגת גבינה & 26 & 23.4 \\
לימונדה & 14 & 12.6 \\
קפה & 12 & 10.8 \\
\bottomrule
\end{tabular}
\end{outtable}

{\small \textbf{הנחה של \lr{10\%}} פירושה שנשאר \lr{90\%}, ולכן מכפילים ב-\sqlin{0.9}. גם \sqlin{DishPrice - DishPrice * 0.1} נכון. \textbf{טעות נפוצה:} להכפיל ב-\sqlin{0.1} \textendash{} זה מחשב את ההנחה עצמה, לא את המחיר החדש.}

\newpage
\partheader{פתרונות \textendash{} שאלה 3}{\texttt{cafe.sql}}{20 נקודות}

\subqp{א}{מנות שמחירן מעל \lr{30}}{5}
\begin{otherlanguage}{english}
\begin{sqlverb}
SELECT DishName, DishPrice
FROM   Dishes
WHERE  DishPrice > 30;
\end{sqlverb}
\end{otherlanguage}

\begin{outtable}{ \textendash{} \lr{3} שורות}
\begin{tabular}{ll}
\toprule
\thd{DishName} & \thd{DishPrice} \\
\midrule
סלט ירקות & 32 \\
שניצל & 56 \\
פסטה & 48 \\
\bottomrule
\end{tabular}
\end{outtable}

{\small \textbf{שימו לב ל"מעל":} קנאפה עולה \lr{30} בדיוק, ולכן היא \textbf{אינה} בפלט. אם היה כתוב "\lr{30} ומעלה" היינו כותבים \sqlin{>= 30}, והיא כן הייתה נכנסת.}

\subqp{ב}{הזמנות מאפריל עד יוני}{5}
\begin{otherlanguage}{english}
\begin{sqlverb}
SELECT OrderNum, CustName, OrderDate
FROM   Orders
WHERE  OrderDate BETWEEN '2026-04-01' AND '2026-06-30';
\end{sqlverb}
\end{otherlanguage}

\begin{outtable}{ \textendash{} \lr{6} שורות}
\begin{tabular}{lll}
\toprule
\thd{OrderNum} & \thd{CustName} & \thd{OrderDate} \\
\midrule
105 & ליאת בר & 2026-04-09 \\
106 & ראמי סעיד & 2026-04-27 \\
107 & יעל שגיא & 2026-05-11 \\
108 & חנאן זועבי & 2026-05-30 \\
109 & איתי פלד & 2026-06-14 \\
110 & מאיה לוי & 2026-06-28 \\
\bottomrule
\end{tabular}
\end{outtable}

{\small התאריך בשאלה כתוב \lr{1.4.2026}, אבל ב-\lr{SQL} כותבים אותו תמיד \sqlin{'2026-04-01'} \textendash{} שנה, חודש, יום, ובגרש בודד. \lr{\texttt{BETWEEN}} \textbf{כולל} את שני הקצוות.}

\subqp{ג}{עובדים בלי טלפון}{5}
\begin{otherlanguage}{english}
\begin{sqlverb}
SELECT WorkerNum, WorkerName
FROM   Workers
WHERE  WorkerPhone IS NULL;
\end{sqlverb}
\end{otherlanguage}

\begin{outtable}{ \textendash{} \lr{2} שורות}
\begin{tabular}{ll}
\toprule
\thd{WorkerNum} & \thd{WorkerName} \\
\midrule
2 & יובל נחום \\
5 & רים עבאס \\
\bottomrule
\end{tabular}
\end{outtable}

\begin{notepanel}{warnred}{הטעות שחייבים לא לעשות}
\sqlin{WHERE WorkerPhone = NULL} מחזיר \textbf{אפס שורות} \textendash{} ובלי להודיע על שגיאה. אי אפשר להשוות לערך שאינו ידוע. הדרך היחידה היא \sqlin{IS NULL}, וההפך הוא \sqlin{IS NOT NULL}.
\end{notepanel}

\subqp{ד}{לקוחות בשם משפחה "לוי"}{5}
\begin{otherlanguage}{english}
\begin{sqlverb}
SELECT OrderNum, CustName
FROM   Orders
WHERE  CustName LIKE '%לוי%';
\end{sqlverb}
\end{otherlanguage}

\begin{outtable}{ \textendash{} \lr{2} שורות}
\begin{tabular}{ll}
\toprule
\thd{OrderNum} & \thd{CustName} \\
\midrule
101 & דנה לוי \\
110 & מאיה לוי \\
\bottomrule
\end{tabular}
\end{outtable}

{\small בעמודה שמור \textbf{השם המלא}, ולכן \sqlin{= 'לוי'} יחזיר אפס שורות. "מכיל" דורש \textbf{שני} סימני אחוז, אחד מכל צד.}

\newpage
\partheader{פתרונות \textendash{} שאלה 4}{\texttt{cafe.sql}}{15 נקודות}

\subqp{א}{מיון לפי סוג, ובתוכו לפי מחיר יורד}{5}
\begin{otherlanguage}{english}
\begin{sqlverb}
SELECT   DishName, DishKind, DishPrice
FROM     Dishes
ORDER BY DishKind, DishPrice DESC;
\end{sqlverb}
\end{otherlanguage}

\begin{outtable}{ \textendash{} \lr{8} שורות}
\begin{tabular}{lll}
\toprule
\thd{DishName} & \thd{DishKind} & \thd{DishPrice} \\
\midrule
שניצל & מנה עיקרית & 56 \\
פסטה & מנה עיקרית & 48 \\
סלט ירקות & מנה ראשונה & 32 \\
חומוס & מנה ראשונה & 28 \\
קנאפה & קינוח & 30 \\
עוגת גבינה & קינוח & 26 \\
לימונדה & שתייה & 14 \\
קפה & שתייה & 12 \\
\bottomrule
\end{tabular}
\end{outtable}

\begin{notepanel}{tipgold}{למה "מנה עיקרית" לפני "מנה ראשונה"?}
כי המיון הוא \textbf{אלפביתי}, ולא לפי סדר ההגשה בתפריט. שתי המחרוזות מתחילות ב"מנה~", וההכרעה נופלת על האות הבאה: \textbf{ע} באה לפני \textbf{ר} באלף-בית. המחשב אינו יודע מה זה "מנה ראשונה" \textendash{} הוא רק משווה אותיות.

\vspace{3pt}
\textbf{וגם:} ה-\lr{\texttt{DESC}} חל \textbf{רק} על \lr{\texttt{DishPrice}}. אילו היה כתוב \sqlin{ORDER BY DishKind DESC, DishPrice} התוצאה הייתה שונה לגמרי.
\end{notepanel}

\subqp{ב}{שלוש ההזמנות עם הכי הרבה סועדים}{5}
\begin{otherlanguage}{english}
\begin{sqlverb}
SELECT   OrderNum, CustName, Guests
FROM     Orders
ORDER BY Guests DESC
LIMIT    3;
\end{sqlverb}
\end{otherlanguage}

\begin{outtable}{ \textendash{} \lr{3} שורות}
\begin{tabular}{lll}
\toprule
\thd{OrderNum} & \thd{CustName} & \thd{Guests} \\
\midrule
112 & שני אדרי & 10 \\
108 & חנאן זועבי & 8 \\
104 & סמיר דיב & 6 \\
\bottomrule
\end{tabular}
\end{outtable}

{\small "הכי הרבה" \textendash{} \lr{\texttt{DESC}}. "שלוש" \textendash{} \sqlin{LIMIT 3}. \textbf{סדר הכתיבה חשוב:} \lr{\texttt{LIMIT}} תמיד \textbf{אחרי} \lr{\texttt{ORDER BY}}, ובלי מיון הוא יחזיר שלוש שורות שרירותיות.}

\subqp{ג}{עובדים לפי שכר יורד, ובשכר זהה לפי שם}{5}
\begin{otherlanguage}{english}
\begin{sqlverb}
SELECT   WorkerName, HourPay
FROM     Workers
ORDER BY HourPay DESC, WorkerName;
\end{sqlverb}
\end{otherlanguage}

\begin{outtable}{ \textendash{} \lr{6} שורות}
\begin{tabular}{ll}
\toprule
\thd{WorkerName} & \thd{HourPay} \\
\midrule
רים עבאס & 55 \\
אמיר טאהא & 50 \\
סאלי חדאד & 42 \\
יובל נחום & 38 \\
תום בראל & 38 \\
מיכל שרון & 36 \\
\bottomrule
\end{tabular}
\end{outtable}

{\small \textbf{הבדיקה המהירה:} שני העובדים שמשתכרים \lr{38} חייבים להופיע בסדר \textbf{יובל ואז תום} \textendash{} י' לפני ת' באלף-בית. אם קיבלתם אותם בסדר הפוך, המפתח השני חסר בשאילתה שלכם.}

\vspace{8pt}
{\color{inkblue}\hrule height 1pt}
\vspace{5pt}
\begin{panel}
\textbf{בדקו את עצמכם לפני בחינת האמצע}

\vspace{3pt}
\noindent\begin{tabular}{@{}p{0.16\linewidth}p{0.78\linewidth}@{}}
\toprule
\thd{הציון} & \thd{מה לעשות עכשיו} \\
\midrule
\lr{50\textendash60} & אתם מוכנים. עברו שוב על דף העזר, ובמיוחד על \lr{\texttt{NULL}} ועל הסוגריים \\
\lr{40\textendash49} & כמעט. סמנו את הסעיפים שטעיתם בהם, וחיזרו על הפרק שלהם בחומר הלימוד \\
\lr{מתחת ל-40} & חיזרו על חומר הלימוד מההתחלה, ופתרו שוב את המיני-תרגול שבסופו \\
\bottomrule
\end{tabular}

\vspace{5pt}
{\small\textbf{שלוש הטעויות שתופסות את רוב הכיתה:} \sqlin{= NULL} במקום \sqlin{IS NULL} \textbullet{} \lr{\texttt{LIKE}} בלי סימני אחוז \textbullet{} \lr{\texttt{DESC}} שנכתב רק ליד העמודה הראשונה.}
\end{panel}

\end{document}
